% LaTeX version of paper/experiment_setting/experiment_setting.md
% Intended to be included via \input{} in a larger manuscript.

\subsection{Experimental Testbed}
\label{sec:experimental-testbed}

We evaluate online budget estimation on four complementary environments that stress distinct modes of agentic decision-making: irreversible planning in Sokoban, search-grounded question answering in SearchR1, repository-level software engineering in a SWE-bench-style setting, and multi-resource supply-chain control in Warehouse. The first three benchmarks target token-budget awareness. Given a partial rollout, the evaluator must decide whether the agent can still finish successfully within a fixed total token budget and, if so, predict an interval for the additional tokens required from the next turn onward. Warehouse targets financial-budget awareness. Given a partial rollout, the evaluator must decide whether the trajectory can still reach a target cash threshold while satisfying joint budgets on time, cumulative warehouse occupancy, and cumulative cost, and, if so, predict one interval for each remaining resource.

We choose these four environments because they expose qualitatively different budget-estimation failure modes. Sokoban emphasizes long-horizon commitment under effectively irreversible actions; SearchR1 emphasizes adaptive search and information synthesis; SWE-bench-style coding emphasizes iterative debugging, file inspection, and repair loops with heavy tool use; and Warehouse emphasizes external resource accumulation in a realistic control setting. These environments also span different interaction horizons. For SearchR1 and Sokoban, the realized episode length is determined by the benchmark-specific \texttt{max\_context\_window\_tokens} budget rather than a fixed 10-turn cap: SearchR1 uses a 3{,}500-token context window with one action per turn, Sokoban uses a 2{,}500-token context window with up to three actions per turn, and the current SWE-bench-style coding runs produce substantially longer and more variable traces, with family-level average trajectory lengths ranging from 12.1 turns (Claude Opus 4.7 Low Thinking) to 62.9 turns (Qwen3 235B). Together, they let us test whether budget awareness is a robust capability across planning, search, coding, and operations domains rather than an artifact of a single benchmark family.

For the interactive language-agent benchmarks, we consider five representative frontier model families: GPT-5.2 Instant, Claude Opus 4.7 Low Thinking, Claude Sonnet 4.6 Low Thinking, Gemini 3.1 Pro Preview, and Qwen3 235B. This set spans multiple providers, both proprietary and open model families, and different reasoning styles, allowing us to test whether budget estimation behavior is model-specific or systematic across frontier agents.

All estimation experiments are conducted on logged rollouts rather than newly regenerated trajectories. SearchR1, Sokoban, and Warehouse contain 128 rollout episodes per benchmark-model configuration when available. The current SWE-bench coding benchmark uses recorded origin trajectories from five rollout families, with 64 rollouts each for GPT-5.2 Instant, Claude Opus 4.7 Low Thinking, Claude Sonnet 4.6 Low Thinking, and Gemini 3.1 Pro Preview, and 60 rollouts for Qwen3 235B. From every non-terminal rollout prefix, we construct one candidate estimation sample by replaying the observed dialogue history and appending a task-specific estimation question. As a result, the number of candidate estimation samples depends on realized trajectory length rather than being fixed in advance.

\section{Detailed Experimental Settings}
\label{sec:detailed-experimental-settings}

\subsection{Environments and Tasks}
\label{sec:environments-and-tasks}

Our benchmark spans four environments with complementary cost structures and action spaces. SearchR1, Sokoban, and SWE-bench use token budgets, whereas Warehouse uses a joint financial budget defined over multiple resources.

\paragraph{Sokoban.}
We use a coordinate-based Sokoban environment to study budget estimation under irreversible planning. The rollout agent operates on procedurally generated puzzles with an $8 \times 8$ grid, two boxes, and search depth 30. The agent operates under a token budget of 2{,}500. A single incorrect push can permanently destroy solvability, so the benchmark stresses whether the agent can recognize when a trajectory is still likely to finish efficiently and when it is already drifting toward failure. For rollout generation we cap each response at 800 tokens and truncate dialogue history at 2{,}500 context tokens.

\paragraph{SearchR1.}
We use the SearchQA/SearchR1 environment backed by a HotpotQA-derived parquet dataset and a retrieval server. At each turn the agent must choose exactly one action, either issuing a search query or submitting a final answer. The agent operates under a token budget of 3{,}500. This benchmark stresses online budget estimation under open-domain information gathering, where token usage depends on how efficiently the agent searches, reformulates queries, and decides when evidence is sufficient to answer.

\paragraph{SWE-bench-style coding.}
We evaluate token estimation on repository-level software-engineering trajectories converted from recorded coding-agent rollouts stored as raw \texttt{*.traj.json} execution logs together with resolution summaries. The current main setting includes five rollout families: GPT-5.2 Instant, Claude Opus 4.7 Low Thinking, Claude Sonnet 4.6 Low Thinking, Gemini 3.1 Pro Preview, and Qwen3 235B. Each converted trajectory is represented as an alternating sequence of user-visible repository context, assistant actions, and subsequent execution feedback, so one turn corresponds to one assistant step together with the resulting tool outputs or validation results. Rollout success is defined from the recorded SWE-bench resolution summaries rather than from patch submission alone. We include this benchmark because token usage in coding agents is dominated by repeated inspection, targeted edits, test execution, and repair iterations, making it a useful stress test for budget estimation in tool-rich, nontrivial workflows.

\paragraph{Warehouse.}
We use logged rollouts from an Acme consumer-electronics supply-chain simulator. Each turn advances 14 days. The rollout agent can place production orders for five SKUs with air or ocean shipping, replenish five retailers by truck, draw or repay credit, factor accounts receivable, or pass. The state reports cash, accounts receivable, credit usage, warehouse and retailer inventory, in-transit and in-production goods, recent demand, events, cash flow, and period P\&L. The original rollout objective is to maximize cumulative profit while avoiding stockouts and bankruptcy, while our estimation task converts each logged prefix into a target-cash feasibility problem under three external resource budgets. The current Warehouse rollout files contain 128 rollouts per model family, with 11 logged assistant turns per rollout and a final state at Day 154, i.e., 22 weeks. The detailed budget construction for this environment is described in Section~\ref{sec:warehouse-budget-construction}.

\subsection{Evaluation Protocol}
\label{sec:evaluation-protocol}

At each evaluation point, we replay the rollout prefix as a dialogue history and ask a separate estimator model to reason about the remaining budget required from the next turn onward. For token-budget benchmarks, the estimator receives the completed rollout history together with a structured summary of per-turn token usage. The estimator must output either an interval $[est_{\mathrm{low}}, est_{\mathrm{high}}]$ for the remaining total tokens or \texttt{impossible}.

Our goal is not to determine whether the environment state admits some successful continuation in principle, but whether the ongoing agent trajectory is still likely to finish successfully under the remaining budget. Throughout this paper, we therefore use feasibility to mean trajectory-level realized feasibility unless explicitly stated otherwise.

Estimator and rollout roles are separated in all experiments. For each benchmark-model entry, we evaluate one designated estimator model $M_{\mathrm{est}}$ on logged rollouts from a specified rollout source. In the main experiments, $M_{\mathrm{est}}$ is the same model family or variant as the rollout generator (self-estimation). Cross-model estimation is supported by pairing a different $M_{\mathrm{est}}$ with the same logged rollouts and is reported separately in the appendix when used. The estimator only receives prefix-available information: replayed history up to the current turn, benchmark budget limits, and structured usage summaries from completed turns. It is not given any future turns, future tool outputs, terminal success labels, or trajectory-level realized remaining budget. We do not provide rollout-model identity as a separate structured input field.

For token tasks, the usage summaries follow the \texttt{turn\_excluding\_history} convention, i.e., per-turn fresh token usage excluding replayed history. The target is future incremental token cost, not prompt replay cost. For Warehouse, the estimator receives target cash and the three resource budgets, \texttt{time\_weeks}, \texttt{warehouse\_item\_weeks}, and \texttt{cumulative\_cost\_usd}, together with cumulative usage so far. Outputs are parsed with strict formats: token tasks require \verb|<answer>[est_low, est_high]</answer>| or \verb|<answer>impossible</answer>|, and Warehouse requires three labeled intervals or \texttt{impossible}. Malformed outputs are treated as missing predictions.

For Warehouse, we use an analogous protocol but with multi-resource financial budgets rather than token budgets. The estimator receives the rollout prefix together with summaries of completed progress and resource usage, and must predict whether the trajectory can still finish while satisfying the time, warehouse-occupancy, and cumulative-cost constraints and reaching the target cash threshold; if so, it outputs one interval for each remaining resource. The detailed budget construction is described in Section~\ref{sec:warehouse-budget-construction}.

Trajectory-level realized feasibility is defined from the full logged rollout. For token-budget tasks, a prefix is feasible if and only if the rollout eventually succeeds and the total token usage of the complete rollout stays within the benchmark-specific budget. We use a 3{,}500-token budget for SearchR1 and a 2{,}500-token budget for Sokoban. For the current SWE-bench coding benchmark, we do not impose a single shared token limit across all model families. Instead, for each rollout family we set \texttt{max\_context\_window\_tokens} to the median realized total rollout tokens under the \texttt{turn\_excluding\_history} accounting used by the estimator. This yields family-specific budgets of 8{,}265 tokens for GPT-5.2 Instant, 7{,}060 for Claude Opus 4.7 Low Thinking, 6{,}772 for Claude Sonnet 4.6 Low Thinking, 10{,}472 for Gemini 3.1 Pro Preview, and 21{,}711 for Qwen3 235B. This choice keeps the coding benchmark comparable across markedly different trajectory-length regimes without forcing all rollout families into a single arbitrary token ceiling.

Conceptually, each non-terminal rollout prefix yields one evaluation sample, excluding the terminal prefix because no future budget remains to be estimated. Thus, if a trajectory has $T$ turns, it contributes $T-1$ candidate estimation samples. In the current SWE-bench runs, however, the flattened coding benchmark produces between 712 and 3{,}715 candidate prefixes per rollout family because the trajectories are much longer than in the other domains. We therefore construct a capped 512-prefix estimation split separately for each rollout family. Let $\mathcal{R}$ denote the rollouts in one family, and let $T_r$ be the number of assistant turns in rollout $r$. We first sort rollouts by $(T_r,r)$ and partition the sorted list into eight nearly equal-cardinality rollout-length buckets, with any remainder assigned to the earlier buckets. This buckets trajectories by realized length at the rollout level rather than by the number of generated prefixes, so long rollouts do not receive proportionally more sampling mass simply because they contain more non-terminal prefixes. Within each rollout, its candidate prefix indices are shuffled using a pseudorandom generator with seed 42; the rollout order inside each length bucket is also shuffled with the same seeded generator. Samples are then selected by repeated randomized round-robin passes: at each pass, the active buckets are randomly ordered, one bucket is visited at a time, and that bucket contributes the next available prefix from the next rollout in its internal round-robin order. Buckets and rollouts with no remaining candidate prefixes are removed, and the procedure stops when 512 prefixes have been selected or all candidates are exhausted. The selected prefixes are finally reordered only for execution efficiency, grouping prefixes from the same rollout and sorting them by observed prefix token usage and turn index to improve prompt-cache reuse; this cache-friendly ordering does not change the membership of the 512-prefix split. This procedure preserves the online prefix-estimation problem while balancing coverage across short and long coding trajectories.

We report two classes of metrics. First, feasibility accuracy measures whether the estimator correctly predicts whether the rollout can still finish within budget. Second, interval quality measures whether the predicted interval contains the trajectory-level realized remaining budget on feasible prefixes. For token-budget benchmarks, we additionally track interval width and a simple aggregate reward that averages feasibility correctness with interval hit on trajectory-level realized feasible prefixes. In the main analysis, we also aggregate samples by relative rollout progress to study how estimation quality changes from early to late stages of execution.

\subsection{Warehouse Budget Construction}
\label{sec:warehouse-budget-construction}

Warehouse is evaluated as a multi-resource financial estimation task rather than a token-budget task. For each prefix, the estimator receives the logged rollout history together with structured summaries of completed turns, completed weeks, current cash, target cash, cumulative time usage, cumulative warehouse occupancy in item-weeks, cumulative cost in USD, and per-step resource consumption observed so far. It must decide whether the logged trajectory can still finish successfully while staying within all three budgets and reaching the target final cash threshold. Here, \texttt{time\_weeks} limits total elapsed time over the full trajectory, \texttt{warehouse\_item\_weeks} limits cumulative storage usage over time, defined as inventory level integrated across weeks, and \texttt{cumulative\_cost\_usd} limits cumulative operational spending. Revenue does not reduce \texttt{cumulative\_cost\_usd}. If the trajectory remains feasible, the estimator must output one interval each for the remaining \texttt{time\_weeks}, \texttt{warehouse\_item\_weeks}, and \texttt{cumulative\_cost\_usd}; otherwise it must answer \texttt{impossible}.

The underlying Warehouse simulator models a consumer-electronics brand managing five SKUs across an Acme warehouse and five retailers. On each 14-day turn, the agent may order production by SKU, order size, and shipping mode; replenish retailers; draw or repay a credit line; factor accounts receivable; or pass. Production, shipping, retailer demand, Net-30 payment delays, holding cost, fixed opex, credit interest, stockout penalties, and overstock penalties all affect the logged state and final cash. This creates coupled trade-offs: faster air production can improve sales but raises cost, replenishment changes both warehouse occupancy and downstream stockouts, and credit or factoring can improve short-term cash while increasing later financing pressure.

The trajectory-level realized label is derived from the full logged rollout. A prefix is feasible only if the logged rollout succeeds, final cash reaches or exceeds the target cash, realized total time is within budget, realized cumulative warehouse occupancy is within budget, and realized cumulative cost is within budget. Realized total time is parsed from the final state's day count; in the current logged Warehouse data all rollouts end at Day 154, giving 22 total weeks. Realized warehouse occupancy and cost are taken from the final turn's \texttt{financials.cumulative\_inventory\_weeks} and \texttt{financials.cumulative\_cost}. Prefix-level remaining labels subtract the current cumulative values from these final realized totals.

In the current main setting we use the \texttt{half\_reachable} preset with random seed 42. The script defaults are \texttt{AUTO\_TIME\_BUDGET\_RATIO=1.0}, \texttt{AUTO\_WAREHOUSE\_BUDGET\_RATIO=1.2}, and \texttt{AUTO\_COST\_BUDGET\_RATIO=1.2}. After shuffling rollouts, half are assigned reachable constraints and half unreachable constraints. For the reachable half, target cash is sampled as $U(0.50,1.00)\times$ final cash and clipped so that it is no greater than final cash; the time budget is exactly $1.0\times$ realized total time; and the warehouse and cost budgets are sampled uniformly between $1.0\times$ and $1.2\times$ their realized final totals. For the unreachable half, target-cash, time-budget, warehouse-budget, and cost-budget ratios are independently resampled from $U(0.50,2.00)$ until the logged rollout fails the feasibility check. Thus unreachable examples can be caused by the target cash, one or more resource budgets, or their combination.

We report feasibility accuracy, per-dimension coverage for time, warehouse occupancy, and cost, joint coverage when all three intervals contain the trajectory-level realized remaining values, midpoint absolute error for each dimension, and an aggregate reward. The reward always includes the feasibility judgment term; if the prefix is trajectory-level realized feasible and the model predicts feasibility, we additionally average in the three interval-coverage terms. This gives a compact scalar score while still preserving per-dimension diagnostics.

\subsection{Experiment Matrix}
\label{sec:experiment-matrix}

\begin{table*}[t]
\centering
\small
\setlength{\tabcolsep}{2pt}
\renewcommand{\arraystretch}{1.1}
\begin{tabular}{|p{0.11\textwidth}|p{0.12\textwidth}|p{0.24\textwidth}|p{0.06\textwidth}|p{0.20\textwidth}|p{0.15\textwidth}|}
\hline
Benchmark & Budget Type & Rollout Models / Sources & Rollouts & Horizon & Estimation Target \\
\hline
SearchR1 & Token budget & GPT-5.2 Instant, Claude Opus 4.7 Low Thinking, Claude Sonnet 4.6 Low Thinking, Gemini 3.1 Pro Preview, Qwen3 235B & 128 & Realized turns depend on \texttt{max\_context\_window\_tokens} $= 3{,}500$; 1 action per turn & Remaining total tokens within 3{,}500 \\
\hline
Sokoban & Token budget & GPT-5.2 Instant, Claude Opus 4.7 Low Thinking, Claude Sonnet 4.6 Low Thinking, Gemini 3.1 Pro Preview, Qwen3 235B & 128 & Realized turns depend on \texttt{max\_context\_window\_tokens} $= 2{,}500$; up to 3 actions per turn & Remaining total tokens within 2{,}500 \\
\hline
SWE-bench-style coding & Token budget & GPT-5.2 Instant, Claude Opus 4.7 Low Thinking, Claude Sonnet 4.6 Low Thinking, Gemini 3.1 Pro Preview, Qwen3 235B (recorded origin coding-agent rollouts) & 64 / 64 / 64 / 64 / 60 & Multi-turn coding trajectories converted from \texttt{*.traj.json} logs; average turns range from 12.1 to 62.9 across rollout families & Remaining total tokens within a rollout-family-specific median budget (6{,}772--21{,}711); 512 sampled prefixes per family \\
\hline
Warehouse & Financial budget (time, warehouse occupancy, cost) & GPT-5.2 Instant, Claude Opus 4.7 Low Thinking, Claude Sonnet 4.6 Low Thinking, Gemini 3.1 Pro Preview, Qwen3 235B & 128 & 11 logged turns; 2 weeks per turn; final Day 154 & Remaining time, warehouse item-weeks, and cost while reaching target cash \\
\hline
\end{tabular}
\caption{Experiment matrix for the four evaluation benchmarks.}
\label{tab:experiment-matrix}
\end{table*}

\subsection{Model Identifiers and Inference Settings}
\label{sec:model-identifiers-and-inference-settings}

To ensure reproducibility, we report exact deployment metadata for every model used in rollout generation and estimation. The display names in the main text are shorthand labels; the appendix table provides the exact API model identifier and invocation configuration used in experiments.

For each model, we log the following fields: display name, provider, exact API model id, or checkpoint hash for local models, date queried, endpoint or region, reasoning configuration, max output tokens, temperature, top-$p$ if used, stop settings, and sampling seed policy. We also record whether retries or fallback routes were enabled.

The term ``Low Thinking'' denotes a constrained reasoning configuration in provider-native settings, not a prompt-only alias. Concretely, this corresponds to provider-specific controls such as \texttt{reasoning\_effort=low}, \texttt{output\_config.effort=low}, or \texttt{thinking\_mode=low}, depending on the API backend. We report the exact control knob and value per model in the appendix table.

Table~\ref{tab:model-identifiers} lists the concrete entries used in this paper, extracted from \texttt{config/evaluate\_api\_llm.yaml}.

\begin{table*}[t]
\centering
\scriptsize
\setlength{\tabcolsep}{1pt}
\renewcommand{\arraystretch}{1.1}
\begin{tabular}{|p{0.105\textwidth}|p{0.09\textwidth}|p{0.055\textwidth}|p{0.12\textwidth}|p{0.095\textwidth}|p{0.06\textwidth}|p{0.045\textwidth}|p{0.08\textwidth}|p{0.08\textwidth}|p{0.065\textwidth}|p{0.085\textwidth}|}
\hline
Display Name in Paper & YAML Key & Provider & API / Model Identifier & Reasoning Setting & Max Output Tokens & Temperature & Date Queried & Endpoint / Region & Retry Policy & Seed Policy \\
\hline
GPT-5.2 Instant & OpenAI-5.2-Instant & openai & gpt-5.2 & reasoning\_effort: none & 800 (max\_completion\_tokens) & N/A & Logged per run in metadata & Logged per run in metadata & Logged per run in metadata & No fixed API sampling seed; deterministic decoding settings reported \\
\hline
Claude Opus 4.7 Low Thinking & Claude-Opus-4.7-low-thinking & anthropic & claude-opus-4-7 & output\_config.effort: low & 800 (max\_tokens) & N/A & Logged per run in metadata & Logged per run in metadata & Logged per run in metadata & No fixed API sampling seed; deterministic decoding settings reported \\
\hline
Claude Sonnet 4.6 Low Thinking & Claude-Sonnet-4.6-low-thinking & anthropic & claude-sonnet-4-6 & output\_config.effort: low & 800 (max\_tokens) & N/A & Logged per run in metadata & Logged per run in metadata & Logged per run in metadata & No fixed API sampling seed; deterministic decoding settings reported \\
\hline
Gemini 3.1 Pro Preview & OpenRouter-Gemini-3.1-Pro-Preview & openrouter & google/gemini-3.1-pro-preview & thinking\_mode: low & 800 (max\_tokens) & 0 & Logged per run in metadata & Logged per run in metadata & Logged per run in metadata & No fixed API sampling seed; deterministic decoding settings reported \\
\hline
Qwen3 235B & qwen/qwen3-235b-a22b-2507 & openrouter & qwen/qwen3-235b-a22b-2507 & thinking\_mode: low & 800 (max\_tokens) & 0 & Logged per run in metadata & Logged per run in metadata & Logged per run in metadata & No fixed API sampling seed; deterministic decoding settings reported \\
\hline
\end{tabular}
\caption{Model deployment identifiers and inference settings used for rollout generation and estimation.}
\label{tab:model-identifiers}
\end{table*}

All fields above are versioned together with evaluation outputs. In the release package, we provide the full run metadata, including query timestamps, routing endpoint information, retry records, and decoding and sampling configuration, alongside the prediction files.
